\section{Setup}

The environment is a finite belief-elicitation problem.
There are \(k\) possible outcomes, and the subject has latent beliefs
\[
p=(p_1,\ldots,p_k)\in\DeltaK,
\quad
p_i\ge 0,
\quad
\sum_{i=1}^k p_i=1.
\]
The elicitation task asks about \(n\in\N\) independent repetitions.
If the subject's beliefs are \(p\), the realized count vector is
\[
\omega=(\omega_1,\ldots,\omega_k)\sim \mathrm{Mult}(n,p).
\]
The set of feasible count vectors is
\[
\OmegaN=\left\{\omega\in\N_0^k:\sum_{i=1}^k \omega_i=n\right\},
\]
and the multinomial probability mass function is
\[
\Pr_p(\omega)
=
\frac{n!}{\omega_1!\cdots \omega_k!}
\prod_{i=1}^k p_i^{\omega_i}.
\]
The subject reports another feasible count vector
\[
r=(r_1,\ldots,r_k)\in\OmegaN.
\]

\begin{definition}[Frequency-report scoring rule]
A frequency-report scoring rule is a function
\[
S:\OmegaN\times\OmegaN\to\R,
\]
where \(S(r,\omega)\) is the score assigned to report \(r\) when the realized count vector is \(\omega\).
The paper studies rules that reward reports close to realized counts.
A convenient form is
\[
S_D(r,\omega)=a-bD(r,\omega),
\quad b>0,
\]
where \(D:\OmegaN\times\OmegaN\to\R_+\) is a count-error index.
The function \(D\) need not be a metric: squared distance is a divergence-style count loss, not a metric, because it fails the triangle inequality.
\end{definition}

The constants \(a\) and \(b\) do not matter for expected-score rankings as long as \(b>0\).
They matter for implementation.
If scores are paid directly as money and nonnegative payments are desired, one can choose
\[
a\ge b\max_{r,\omega\in\OmegaN}D(r,\omega).
\]
How the score is implemented as a payment affects robustness to the subject's risk attitude; the paper's discussion of risk aversion and implementation takes up this question separately.

\begin{definition}[Optimal reports]
Under risk neutrality, the optimal-report correspondence induced by a scoring rule \(S\) is
\[
R_S(p)
=
\argmax_{r\in\OmegaN}\E_p[S(r,\omega)].
\]
Because \(\OmegaN\) is finite, \(R_S(p)\) is nonempty.
\end{definition}

Risk neutrality is the maintained assumption throughout the body of the paper; the discussion of risk aversion and implementation separately shows how the analysis extends to risk-averse expected-utility subjects.

This is the forward incentive object.
It answers the question: if the subject's latent beliefs are \(p\), which reports are optimal?
The econometric object is the inverse question.
After observing \(r\), the researcher asks which beliefs could have generated that report.

\begin{definition}[Identified set]
For an observed report \(r\in\OmegaN\), the identified set induced by \(S\) is
\[
P_S(r)
=
\{p\in\DeltaK:r\in R_S(p)\}.
\]
Ties are included: if \(r\) is one of several optimal reports under \(p\), then \(p\in P_S(r)\).
\end{definition}

We call \(P_S(r)\) the \emph{mechanism-induced identified set}: the identifying restriction is the optimal-report condition \(r\in R_S(p)\) imposed by the scoring-rule mechanism, not a sampling distribution over repeated observations. By construction \(P_S(r)\) is exactly the set of beliefs consistent with that restriction.

The most direct summaries of \(P_S(r)\) are coordinate-wise identification bounds:
\[
\underline p_{S,i}(r)
=
\inf_{p\in P_S(r)}p_i,
\quad
\overline p_{S,i}(r)
=
\sup_{p\in P_S(r)}p_i,
\]
with interval
\[
P_{S,i}(r)
=
[\underline p_{S,i}(r),\overline p_{S,i}(r)].
\]
When the scoring rule is clear from context, the subscript \(S\) is omitted.

Many experiments require more than coordinate-by-coordinate probability bounds.
If outcomes have numerical values \(x_1,\ldots,x_k\), the latent mean implied by beliefs \(p\) is
\[
\mu(p;x)=\sum_{i=1}^k p_i x_i.
\]
This could represent an expected outcome, an expected treatment share, a predicted group average, or another linear index of the belief distribution.
The sharp mean bounds after observing \(r\) are
\[
\underline\mu_S(r;x)
=
\inf_{p\in P_S(r)}\sum_{i=1}^k p_i x_i,
\quad
\overline\mu_S(r;x)
=
\sup_{p\in P_S(r)}\sum_{i=1}^k p_i x_i.
\]
The corresponding mean interval is
\[
[\underline\mu_S(r;x),\overline\mu_S(r;x)].
\]
These bounds generally require optimizing over the full joint region \(P_S(r)\).
Coordinate intervals alone do not generally give sharp bounds for a mean, because they ignore dependence across coordinates imposed by the simplex and the scoring rule.

The design comparison later summarizes the informativeness of a rule using coordinate widths,
\[
\overline p_{S,i}(r)-\underline p_{S,i}(r),
\]
and functional widths,
\[
\overline\mu_S(r;x)-\underline\mu_S(r;x).
\]
Exact-match payment probability is reported only for the discrete-metric rule, as an implementation diagnostic of that fixed-prize mechanism; it is not a cross-rule comparison metric.

The next section applies this setup to the three headline rules: quadratic-distance, discrete-metric, and Manhattan.
